\chapter{Tracheostomy}

\section*{Introduction \& Clinical Indications}
A tracheostomy is a surgical procedure that involves creating an opening (stoma) in the anterior wall of the trachea, typically at the level of the second or third tracheal ring. This intervention is performed to establish a direct airway that bypasses the upper respiratory tract, facilitating ventilation and airway management in patients with compromised respiratory function. It is particularly significant in cases where prolonged mechanical ventilation is necessary, severe upper airway obstruction is present, or when patients are unable to effectively clear their own airways due to various medical conditions. The tracheostomy tube allows for improved pulmonary hygiene, potential for speech with specialized valves, and often facilitates earlier patient mobilization and rehabilitation. This procedure is a critical component of airway management in both emergency and elective settings, providing a more stable and comfortable long-term airway solution compared to endotracheal intubation.

\begin{itemize}
  \item Tracheotomy
  \item Trach tube insertion
  \item Surgical airway
  \item Percutaneous tracheostomy
  \item Open tracheostomy
  \item Tracheostomy tube placement
\end{itemize}

The primary surgical principle of a tracheostomy is to create a secure and stable airway below the level of the larynx, thereby bypassing any anatomical or functional impairment in the upper airway. This direct access to the trachea serves multiple critical functions. First, it provides a reliable conduit for positive pressure ventilation, which is essential for patients requiring prolonged mechanical ventilation. By reducing the work of breathing, a tracheostomy allows for extended ventilatory support without the complications associated with long-term endotracheal intubation, such as laryngeal injury and tracheal damage.

Additionally, the procedure facilitates effective pulmonary hygiene by allowing for direct suctioning of tracheobronchial secretions, which is vital for patients with impaired cough reflexes or excessive mucus production. Furthermore, a tracheostomy can protect the lower airway from aspiration in patients with severe dysphagia or impaired laryngeal protective reflexes. The technical goals of the procedure include creating a stoma that is appropriately sized and positioned to prevent complications such as tracheal stenosis or erosion into adjacent structures while ensuring a secure fit for the tracheostomy tube.

The history of tracheostomy dates back to ancient civilizations, with early references found in Egyptian and Greek medical texts. Asclepiades of Bithynia, around 100 BC, is often credited with the first recorded tracheostomy, although details are limited. During the Middle Ages, the procedure was rarely performed due to high mortality rates and was often seen as a last resort.

A significant milestone occurred in 1546 when Antonio Musa Brassavola successfully performed a tracheostomy on a patient with a peritonsillar abscess. The 19th century saw advancements in the technique, particularly through the work of Armand Trousseau, who performed numerous tracheostomies during a diphtheria epidemic in Paris, improving survival rates dramatically. However, it was Chevalier Jackson, an American laryngologist, who standardized the modern tracheostomy technique in 1909, emphasizing meticulous surgical technique and postoperative care. His contributions transformed tracheostomy into a safer and more predictable procedure, paving the way for its widespread use in modern critical care. The introduction of percutaneous techniques in the late 20th century further refined the procedure, making it accessible at the bedside.

Tracheostomy is indicated in various clinical scenarios, including:

\begin{itemize}
  \item To provide long-term mechanical ventilatory support for patients requiring ventilation for more than 10-14 days.
  \item To bypass severe upper airway obstruction due to tumors, trauma, or stenosis.
  \item To facilitate the removal of excessive tracheobronchial secretions in patients with impaired cough reflexes.
  \item To protect the airway from aspiration in patients with severe dysphagia or altered mental status.
  \item To reduce the work of breathing in patients with chronic respiratory failure.
  \item As an emergency airway in situations where endotracheal intubation is impossible.
  \item To allow for easier administration of inhaled medications directly into the lower airway.
  \item To improve patient comfort and quality of life in chronic respiratory conditions.
\end{itemize}

Standardized diagnostic scores, such as the Alvarado Score for appendicitis, may be used to assess the need for tracheostomy in specific cases.

Tracheostomy can serve as both a primary and secondary procedure, depending on the clinical context. It is considered a primary procedure when it is the initial surgical intervention chosen for long-term airway management, particularly in cases of anticipated prolonged mechanical ventilation or irreversible upper airway obstruction. For example, a patient with a large laryngeal tumor causing acute obstruction might undergo an immediate tracheostomy.

More commonly, tracheostomy functions as a secondary procedure, following an initial period of endotracheal intubation. In this scenario, it is performed to transition a patient from a temporary translaryngeal airway to a more stable and less damaging long-term solution, typically after 10-14 days of intubation. This secondary role aims to reduce the risks of laryngeal injury and improve patient comfort, facilitating oral feeding and enabling speech with a speaking valve.

\section*{Relevant Surgical Anatomy}
Understanding the relevant surgical anatomy is crucial for performing a tracheostomy safely and effectively. The trachea is a tubular structure that extends from the larynx to the carina, where it bifurcates into the right and left main bronchi. It is composed of C-shaped cartilaginous rings that provide structural support and is lined with respiratory epithelium. The trachea is located anterior to the esophagus and is surrounded by various important structures, including the thyroid gland, major blood vessels, and nerves.

Key anatomical landmarks include:

\begin{itemize}
  \item **Tracheal Rings:** The second and third tracheal rings are typically the sites for incision.
  \item **Thyroid Gland:** The isthmus of the thyroid gland may need to be retracted or divided during the procedure.
  \item **Recurrent Laryngeal Nerve:** This nerve runs in close proximity to the trachea and must be preserved to avoid vocal cord paralysis.
  \item **Major Blood Vessels:** The innominate artery and the carotid arteries are located laterally and must be carefully avoided during dissection.
  \item **Lymphatic Drainage:** The trachea has a rich lymphatic network that drains into the cervical lymph nodes.
\end{itemize}

The surgical approach requires careful identification of these structures to minimize complications and ensure successful placement of the tracheostomy tube.

\section*{Preoperative Workup \& Preparation}
Thorough preoperative preparation is crucial to minimize risks and ensure a smooth procedure.

\begin{itemize}
  \item \textbf{Informed Consent:} Detailed discussion with the patient or their legal proxy regarding the indications, procedure steps, potential benefits, risks, and alternatives.

  \item \textbf{Laboratory Studies:} Routine blood work including a complete blood count (CBC), coagulation profile (PT/INR, PTT), and electrolyte panel.

  \item \textbf{Imaging Studies:} A recent chest X-ray is essential to assess lung fields and rule out pneumothorax. Neck imaging (CT scan) may be useful in complex cases.

  \item \textbf{Anesthesia and Cardiac Clearance:} Comprehensive medical evaluation to assess the patient's fitness for general anesthesia.

  \item \textbf{Medication Review:} Adjustment or temporary discontinuation of anticoagulants or antiplatelet agents as per institutional guidelines.

  \item \textbf{NPO Status:} The patient must be kept nil per os (NPO) for at least 6-8 hours prior to the procedure.

  \item \textbf{Antibiotic Prophylaxis:} Administration of prophylactic broad-spectrum antibiotics, typically within one hour prior to incision.

  \item \textbf{Equipment Preparation:} Ensuring availability of various sizes of tracheostomy tubes, suction equipment, and emergency airway management tools.
\end{itemize}

While tracheostomy is a life-saving procedure, certain conditions may contraindicate or necessitate extreme caution during its performance.

\textbf{Situations requiring an alternative airway plan or modified technique:}
\begin{itemize}
  \item \textbf{Uncorrectable Coagulopathy:} Severe bleeding disorders pose an unacceptably high risk of hemorrhage.
  
  \item \textbf{Severe Anatomical Distortion:} Extreme distortion of the neck due to tumors or trauma that makes identification of tracheal landmarks impossible.

  \item \textbf{Active Infection at the Proposed Surgical Site:} Localized cellulitis or abscess in the neck area can lead to severe wound infection.
\end{itemize}

\textbf{Relative Contraindications and Precautions:}
\begin{itemize}
  \item \textbf{Obesity or Short Neck:} These factors can obscure anatomical landmarks and increase the risk of complications.

  \item \textbf{Previous Neck Surgery or Radiation Therapy:} Scar tissue can increase the difficulty of dissection.

  \item \textbf{Thyroid Gland Enlargement (Goiter):} A large thyroid gland may obstruct access to the trachea.

  \item \textbf{Unstable Cervical Spine Injury:} Neck extension required for the procedure may exacerbate spinal cord injury.

  \item \textbf{High Tracheal Bifurcation or Tracheomalacia:} These conditions can complicate tube placement.

  \item \textbf{Pediatric Patients:} Tracheostomy in children requires specialized techniques due to smaller tracheal rings.
\end{itemize}

\section*{Surgical Technique \& Procedure}
The tracheostomy procedure can be performed as an open surgical tracheostomy or as a percutaneous dilatational tracheostomy (PDT). Both methods aim to create a secure tracheal stoma.

\begin{itemize}
  \item \textbf{Anesthesia and Positioning:} The patient is typically under general anesthesia, often with ongoing endotracheal intubation. The patient is positioned supine with the neck extended using a shoulder roll. The neck is prepped and draped in a sterile fashion.
  
  \item \textbf{Incision (Open Tracheostomy):} A transverse skin incision, approximately 3-4 cm long, is made midway between the cricoid cartilage and the sternal notch. A vertical incision may be used in emergency situations.

  \item \textbf{Dissection:} The platysma muscle is divided, and the strap muscles (sternohyoid and sternothyroid) are identified and retracted laterally. The thyroid isthmus may be retracted or divided to expose the anterior tracheal rings.

  \item \textbf{Tracheal Incision:} The anterior wall of the trachea is identified, usually at the level of the second or third tracheal ring. A small window or H-shaped incision is made, ensuring minimal cartilage removal.

  \item \textbf{Tube Insertion:} A cuffed tracheostomy tube of appropriate size is inserted into the tracheal stoma. The cuff is inflated, and proper placement is confirmed by auscultation of bilateral breath sounds and capnography.

  \item \textbf{Securing the Tube:} The tracheostomy tube is secured to the neck with tracheostomy ties. The skin incision is loosely approximated around the tube to allow for drainage.

  \item \textbf{Percutaneous Dilatational Tracheostomy (PDT):} This technique involves a small skin incision, followed by serial dilation of the pretracheal tissues and tracheal wall over a guidewire, under bronchoscopic guidance.
\end{itemize}

After the procedure, a chest X-ray is often performed to confirm tube position and rule out complications such as pneumothorax.

\section*{Complications \& Risks}
Tracheostomy, like any surgical procedure, carries inherent risks, which can be broadly categorized as general surgical risks and procedure-specific complications.

\textbf{General Surgical Risks:}
\begin{itemize}
  \item \textbf{Bleeding:} Hemorrhage from superficial vessels or major vessels like the innominate artery.

  \item \textbf{Infection:} Surgical site infection or more severe infections like mediastinitis.

  \item \textbf{Anesthesia Complications:} Adverse reactions to anesthetic agents or respiratory depression.
\end{itemize}

\textbf{Procedure-Specific Risks:}
\begin{itemize}
  \item \textbf{Pneumothorax or Pneumomediastinum:} Accidental puncture of the pleura during dissection.

  \item \textbf{Tube Dislodgement or Accidental Decannulation:} This can lead to airway loss and requires immediate re-establishment.

  \item \textbf{Subcutaneous Emphysema:} Air leakage from the trachea into surrounding soft tissues.

  \item \textbf{Tracheal Stenosis:} Narrowing of the trachea at the stoma site, often a late complication.

  \item \textbf{Tracheoesophageal Fistula (TEF):} An abnormal connection between the trachea and esophagus.

  \item \textbf{Tracheoinnominate Artery Fistula (TIF):} A rare but catastrophic complication causing massive hemorrhage.

  \item \textbf{Recurrent Laryngeal Nerve Injury:} Damage to the nerve during dissection.

  \item \textbf{Dysphagia:} Difficulty swallowing, which can be exacerbated by the presence of the tube.

  \item \textbf{Granulation Tissue Formation:} Excessive tissue growth at the stoma site, potentially causing airway obstruction.

  \item \textbf{Cosmetic Scarring:} A visible scar on the neck, which can be a concern for some patients.
\end{itemize}

\section*{Postoperative Care \& Recovery}
Following a tracheostomy, the patient is typically monitored in a critical care unit for close observation. The immediate postoperative focus is on ensuring airway patency, securing tube placement, and managing secretions. The tracheostomy tube cuff is usually inflated to prevent aspiration and facilitate ventilation. 

During the first 24-48 hours, the surgical site is monitored for bleeding, swelling, and signs of infection. Pain management is provided as needed, and patients are gradually mobilized as their condition allows. Within the first week, the tracheostomy tract matures, reducing the risk of accidental decannulation. Patients may begin speech therapy to explore communication options, such as speaking valves, and physical therapy to regain strength and mobility.

Long-term recovery involves ongoing tracheostomy care, including regular cleaning, inner cannula changes, and assessment for decannulation when the underlying condition resolves.

During a tracheostomy, the surgeon documents the condition of the trachea, surrounding structures, and any abnormalities encountered. The pathology report on resected tissues, if applicable, may reveal underlying conditions such as inflammation, infection, or malignancy. The findings can provide insight into the patient's respiratory status and guide further management.

A normal postoperative course following a tracheostomy includes stable oxygenation, effective secretion clearance, and a healing stoma without signs of infection. Patients typically experience pain resolution within a few days, with gradual improvement in respiratory function. The tracheostomy tube should remain patent, and patients should be able to communicate effectively with appropriate interventions, such as speaking valves.

Postoperative care is essential for ensuring optimal recovery and preventing complications. Follow-up care includes:

\begin{itemize}
  \item \textbf{Regular Clinic Visits:} Scheduled appointments to monitor stoma healing and overall patient progress.

  \item \textbf{Tracheostomy Tube Changes:} Initial tube changes are typically performed by healthcare professionals after 7-10 days.

  \item \textbf{Speech and Swallowing Therapy:} Evaluation by a speech-language pathologist to assess swallowing function and introduce speaking valves.

  \item \textbf{Pulmonary Rehabilitation:} Ongoing rehabilitation to optimize lung function and secretion management.

  \item \textbf{Imaging Studies:} Occasional chest X-rays or bronchoscopy to assess tracheal integrity.

  \item \textbf{Patient and Caregiver Education:} Training on tracheostomy care, suctioning techniques, and recognition of warning signs.
\end{itemize}

Patients and caregivers should be educated on warning signs such as difficulty breathing, excessive bleeding, fever, or tube dislodgement, and instructed to seek immediate medical attention if these occur.

\section*{Outcomes \& Prognosis}
Tracheostomy is a common procedure in critical care, with its incidence varying based on patient population and regional practices. The incidence of tracheostomy in mechanically ventilated patients ranges from 5\% to 15\%, with higher rates observed in specialized units. The procedure is performed across all age groups, from neonates to the elderly, with no significant sex predilection. The most common indications for tracheostomy include prolonged mechanical ventilation, upper airway obstruction, and chronic secretion management. While the mortality directly attributable to tracheostomy is low, the overall mortality rate in patients undergoing the procedure reflects the severity of their underlying critical illness, often ranging from 20-50\% in the ICU setting.

The outcome of tracheostomy depends primarily on the indication, neurologic and respiratory function, secretion burden, swallowing, and recovery from the underlying illness. It may facilitate ventilation, pulmonary toilet, communication, or weaning, but does not itself guarantee decannulation or oral feeding. Hemorrhage, obstruction, accidental decannulation, infection, tracheal stenosis, tracheoesophageal fistula, and persistent need for the tube remain possible; decannulation rates vary and should not be represented by one universal percentage.

\section*{References}
\begin{enumerate}
  \item MedlinePlus. Tracheostomy. U.S. National Library of Medicine; 2024. Available from: \url{https://medlineplus.gov/tracheostomy.html}
  
  \item Townsend CM, Beauchamp RD, Evers BM, Mattox KL, editors. Sabiston Textbook of Surgery: The Biological Basis of Modern Surgical Practice. 21st ed. Elsevier; 2022.

  \item American Academy of Otolaryngology--Head and Neck Surgery. Clinical Practice Guideline: Tracheostomy. 2023. Available from: \url{https://www.entnet.org/resource/clinical-practice-guideline-tracheostomy/}

  \item Dulguerov P, Gysin C. Tracheostomy: from antiquity to the third millennium. Laryngoscope. 2007;117(10):1711-1718.
\end{enumerate}

\clearpage
